\begin{abstract}
Feed-forward 3D reconstruction systems routinely resize, crop, and pad their
input images, although these operations change the image formation model.  An
image-space affine transform $A$ changes the camera matrix from $K[R\mid t]$
to $AK[R\mid t]$; after this known update is accounted for, the recovered
scene should remain unchanged up to the model's global gauge.  We formulate
this requirement as a preprocessing-equivariance contract and introduce an
auditable protocol that records both user-level and model-internal image
transforms, composes their exact pixel maps, and compares reconstructed
cameras, intrinsics, depth, and point geometry against ground truth.  Across
13 ETH3D scenes and 104 complete evaluations, independently shifted 75\%
crops increase median rotation error by \SI{4.61}{\degree} for VGGT and
\SI{4.52}{\degree} for DUSt3R relative to paired identity runs; translation
error rises by \SI{9.49}{\degree} and \SI{10.99}{\degree}.  Depth AbsRel
increases by 2.30 and 1.64 percentage points, and point-cloud completeness
error by \SI{0.369}{\meter} and \SI{0.416}{\meter}; all eight corresponding
scene-bootstrap intervals exclude zero.  In contrast, letterbox depth and
point-geometry deltas are statistically indistinguishable from identity for
both models.  Point accuracy degrades decisively for VGGT but not DUSt3R,
preventing an all-metrics claim.  The evidence therefore establishes a
specific multi-model accuracy failure under view-dependent cropping, not
generic fragility to every resize or crop; cross-dataset generality, failure
detection, and repair remain separate empirical claims.
\end{abstract}
